\begin{titlepage}
\thispagestyle{empty}
\begin{tikzpicture}[remember picture, overlay]
  \fill[teacherblue] (current page.north west) rectangle (current page.south east);
  \fill[chalkgold!18] ([xshift=-0.6cm,yshift=-2.0cm]current page.north east) circle (6.6cm);
  \fill[white!7] ([xshift=1.4cm,yshift=2.2cm]current page.south west) circle (5.1cm);

  \draw[chalkgold!75, line width=1.4pt]
    ([xshift=1.1cm,yshift=-1.1cm]current page.north west)
    rectangle
    ([xshift=-1.1cm,yshift=1.1cm]current page.south east);

  \node[
    anchor=north west,
    fill=chalkgold,
    text=teacherblue,
    font=\bfseries\large,
    rounded corners=4pt,
    inner xsep=10pt,
    inner ysep=6pt
  ] at ([xshift=1.6cm,yshift=-1.5cm]current page.north west) {QUYỂN XII};

  \node[
    anchor=north,
    text=white,
    font=\large\itshape
  ] at ([yshift=-1.7cm]current page.north) {Truyện Tích Cảm Hứng Song Ngữ};

  \node[
    anchor=west,
    text=white,
    font=\fontsize{27}{31}\selectfont\bfseries
  ] at ([xshift=1.9cm,yshift=3.6cm]current page.center) {NHỮNG ĐIỀU};

  \node[
    anchor=west,
    text=chalkgold,
    font=\fontsize{27}{31}\selectfont\bfseries
  ] at ([xshift=1.9cm,yshift=1.8cm]current page.center) {THẦY GIÁO};

  \node[
    anchor=west,
    text=white,
    font=\fontsize{27}{31}\selectfont\bfseries
  ] at ([xshift=1.9cm,yshift=0.0cm]current page.center) {NHÌN THẤY};

  \draw[chalkgold, line width=2pt]
    ([xshift=1.9cm,yshift=-1.2cm]current page.center) --
    ([xshift=9.3cm,yshift=-1.2cm]current page.center);

  \node[
    anchor=west,
    text=white,
    font=\large
  ] at ([xshift=1.9cm,yshift=-2.3cm]current page.center) {Một quyển để tự đọc, học ngữ, và lấy chuyện thật mà dạy học trò};

  \node[
    anchor=north west,
    fill=white,
    text=black,
    rounded corners=8pt,
    text width=11.7cm,
    inner sep=14pt,
    align=left,
    font=\normalsize
  ] at ([xshift=1.9cm,yshift=-14.8cm]current page.north west) {%
    \textbf{Tinh thần quyển này:} ngắn, thẳng, thật, đi từ lớp học ra cuộc đời. Không cố văn hoa, chỉ cố nhìn đúng người, đúng nỗi mệt, đúng chỗ đau và đúng điều đáng thương của học trò.\\[6pt]
    \textbf{Cách dùng:} tự đọc để nghĩ, học tiếng Anh qua các câu ngắn, và rút ví dụ giảng kỹ năng sống cho học sinh.\\[6pt]
    \textbf{Điểm tựa:} thô nhưng thật, gần lớp học, gần gia đình, gần những điều người trẻ ít nói ra.
  };

  \node[
    anchor=south,
    text=chalkgold,
    font=\large\itshape,
    align=center
  ] at ([yshift=2.15cm]current page.south) {%
    ``Có những điều học sinh không nói trong bài viết.\\
    Nhưng nó hiện ra ở ánh mắt, ở im lặng, và ở cách các em ngồi trong lớp.''%
  };

  \node[
    anchor=south,
    text=white!85,
    font=\normalsize,
    align=center
  ] at ([yshift=1.1cm]current page.south) {Dành cho người dạy học, người quan sát con người, và người còn muốn hiểu học trò sâu hơn};
\end{tikzpicture}
\end{titlepage}
\clearpage